\documentclass[11pt]{article}
\usepackage[margin=0.85in]{geometry}
\usepackage{amsmath,amssymb,amsthm,hyperref}
\setlength{\emergencystretch}{1.5em}
\newtheorem{theorem}{Theorem}
\DeclareMathOperator{\Rea}{Re}
\title{An explicit lower bound for zeros in a moving Dirichlet window}
\author{GRH transfer supplement}
\date{26 November 2026}
\begin{document}
\maketitle
\noindent This is a classical paper deduction from Dirichlet GRH, not a new kernel-checked
corollary. The added feature is an elementary numerical bound with a free bandwidth,
valid for every window center. The centered asymptotic phenomenon and a sharper centered
second-order bound are already in Zhao's work.
\begin{theorem}
Let $\chi$ have conductor $q>1$. Assume Dirichlet GRH. Let $H\in\mathbb R$,
$\Delta\ge 1$, $T\Delta\ge 1/2$, and count the nontrivial zeros with
$|\operatorname{Im}\rho-H|<T$, including multiplicities. Then
\[
 N_\chi(H;T)\ \ge\
 \frac{\pi}{16\Delta}\left(1-\frac{1}{4\Delta^2T^2}\right)
 \log\frac{q(2+|H|)}{\pi}
 -\frac{32}{\Delta}-\frac{32 e^{\pi\Delta}}{\Delta^2}.
\]
\end{theorem}
\begin{proof}
Use the primitive inducing character; the additional Euler factors of an imprimitive
character have no zeros in the open critical strip. Put $b=\Delta T$ and use the
Fourier convention $\widehat f(v)=\int_{\mathbb R}f(u)e^{-2\pi iuv}\,du$. Define
\[
 k(v)=\frac{\pi}{2}\cos(\pi v)\mathbf 1_{[-1/2,1/2]}(v),\qquad
 g(u)=\frac{\cos(\pi u)}{1-4u^2},\qquad
 f(t)=\left(1-\frac{(t-H)^2}{T^2}\right)g(\Delta(t-H))^2.
\]
The removable singularities in $g$ are filled by continuity. Integration of $k$
gives $g=\check k$, and $k\ge0$ with integral 1 gives $|g|\le1$ on the real line.
Thus $f(t)\le\mathbf 1_{|t-H|<T}$, including the boundary where $f$ is zero.
The function is entire, with quadratic decay in horizontal strips; its
Fourier transform is supported in $[-\Delta,\Delta]$.

Write $K=k*k$. Fourier scaling and differentiation give
\[
 \widehat f(v)=\frac{e^{-2\pi i Hv}}{\Delta}
 \left(K(v/\Delta)+\frac{K''(v/\Delta)}{4\pi^2 b^2}\right).
\]
There are no boundary atoms in $k'$ because $k$ is continuous. Hence
$K''=k'*k'$, $K(0)=\pi^2/8$ and $K''(0)=-\pi^4/8$ by integrating
$\cos^2$ and $\sin^2$. For $0\le v\le\Delta$ put $s=1-v/\Delta$.
The convolution overlap has length $s$. There $k$ is bounded by
$\pi^2/2$ times distance to its relevant endpoint, so
$|K(v/\Delta)|\le \pi^4s^3/24$ and $|K''(v/\Delta)|\le \pi^4s/4$.
Consequently $|\widehat f(v)|\le 7(\Delta-v)/\Delta^2$, using $b\ge1/2$.

The Guinand--Weil formula for the completed primitive nonprincipal function
has conductor contribution $(2\pi)^{-1}\widehat f(0)\log(q/\pi)$,
archimedean contribution $(2\pi)^{-1}\int f(t)\Rea\psi(\sigma+it/2)\,dt$
with $\sigma=1/4$ or $3/4$, and a prime-power term of absolute value at most
\[
 \frac1\pi\sum_{n\le X}\frac{\Lambda(n)}{\sqrt n}
 |\widehat f(\log n/(2\pi))|,\qquad X=e^{2\pi\Delta}.
\]
For the translated test the two prime terms are the character and its conjugate
times the transforms at opposite frequencies; their absolute values agree.
GRH ensures that the left-hand side of the explicit formula is $\sum_\rho f(\operatorname{Im}\rho)$,
which is bounded above by the desired strict-window count.

For the prime term, $\sum_{n\le x}\Lambda(n)\le4x$ follows by
$\sum_{m<n\le2m}\Lambda(n)\le\log\binom{2m}{m}\le 2m\log 2$, summation at dyadic powers,
and rounding up to the next dyadic power. Now split at $Xe^{-j}$ for $j\ge0$:
\[
 \sum_{n\le X}\frac{\Lambda(n)}{\sqrt n}\log(X/n)
 \le \frac{4\sqrt e}{(1-e^{-1/2})^2}\sqrt X.
\]
Indeed the interval $(Xe^{-(j+1)},Xe^{-j}]$ contributes at most
$4\sqrt e\sqrt X(j+1)e^{-j/2}$. Combining with the linear edge bound gives a prime error
at most $32e^{\pi\Delta}/\Delta^2$.

For the archimedean term, the digamma series and a unit-panel comparison give
$|\psi(z)-\log z|\le 1/\Rea z$ when $\Rea z>0$. In detail, compare each
$(z+n)^{-1}$ with its integral on $[n,n+1]$, bound the error by
$\int_n^{n+1}(\Rea z+x)^{-2}\,dx$, and use
$\psi(z)=\lim_{N\to\infty}(\log N-\sum_{0\le n<N}(z+n)^{-1})$.
For $\sigma$ as above, $\tfrac1{12}(2+|t|)\le|\sigma+it/2|\le 2+|t|$.
Writing $u=\Delta(t-H)$, the two quantities $2+|t|$ and $2+|H|$
are within a factor $1+|u|/2$ of one another. Thus
\[
 |\Rea\psi(\sigma+it/2)-\log(2+|H|)|
 \le 7+\log(1+|u|/2).
\]
For $|u|\le1$ the absolute value of $(1-u^2/b^2)g(u)^2$ is at most 5;
for $|u|\ge1$ it is at most $5/(9u^2)$. On the latter region use
$\log(1+|u|/2)\le\log(3/2)+\log|u|$ and
$\int_1^\infty u^{-2}\log u\,du=1$.
The archimedean term is therefore
$(2\pi)^{-1}\log(2+|H|)\widehat f(0)$ with error at most $32/\Delta$.
Substituting the value of $\widehat f(0)$ proves the claim.
\end{proof}

\paragraph{Conductor thresholds and shrinking windows.}
For a fixed $T>0$, take $\Delta=\max(1,1/T)$, so the coefficient
$1-1/(4\Delta^2T^2)$ is at least $3/4$. At $H=0$, it suffices that
\[
 \log(2q/\pi)\ \ge\ \frac{64}{3\pi}
 \left(M\Delta+32+\frac{32e^{\pi\Delta}}{\Delta}\right)
\]
to guarantee $M$ zeros strictly inside the window. For $T\le1$ a
narrower bandwidth $\Delta=(1+T)/(2T)$ may be used in the full theorem
to reduce the exponential scale of the threshold to $e^{\pi/(2T)}$
times factors of polynomial size in $1/T$.

For a moving window let $U=\log(q(2+|H|)/\pi)$ tend to infinity.
Take $\Delta=(\log U-6)/\pi$ and $T=1/(2\Delta)+1/\Delta^2$.
For $\Delta\ge2$, $1-1/(4\Delta^2T^2)\ge 1/\Delta$.
Since $32 e^{-6}<\pi/32$, the theorem gives
$N_\chi(H;T)\ge \pi U/(32\Delta^2)-32/\Delta$, uniformly in $H$.
In particular every such window has a number of zeros tending to infinity, with
$T=\pi/(2\log U)+O((\log U)^{-2})$.

\paragraph{Provenance and scope.}
Bennett--Martin--O'Bryant--Rechnitzer, \href{https://doi.org/10.1090/mcom/3599}{Counting zeros of Dirichlet \(L\)-functions},
already point out that their qualitative Conjecture 1.3 follows from Selberg under
Dirichlet GRH. Their separate Conjecture 1.4 is not addressed.
See \href{https://arxiv.org/abs/1503.00955}{Carneiro--Chandee--Milinovich} for the
explicit formula with general strip-analytic tests, and
\href{https://arxiv.org/abs/2503.15832}{Zhao, Theorems 1 and 8}, for the related
cosine test, the earlier centered shrinking-window result, and its sharper
second-order constant. No novelty claim is made for the centered result.
The \href{https://math.mit.edu/~drew/shortgrh/}{exact-revision audit}
covers both editions at \texttt{72fead10}, not this deduction.
The latest incoming transfer archive inspected
(\href{https://github.com/VladimirReshetnikov/ProveIt/commit/b35eb211b859d2883c98ef71fbe47e1fcbcc4a74}{\texttt{b35eb21}})
covers different consequences and a cubic-field target.
\end{document}

